% ============================================================
%  H. Høffding, "Filosofien og Darwinismen"
%  Nær og Fjern, 1874
%  TRANSLATION (English)
% ============================================================
\documentclass[12pt,a4paper]{article}
\usepackage[utf8]{inputenc}
\usepackage[T1]{fontenc}
\usepackage[english]{babel}
\usepackage{microtype}
\usepackage{setspace}
\usepackage[margin=1.2in]{geometry}
\usepackage{fancyhdr}
\usepackage{hyperref}

\hypersetup{
  pdftitle={Philosophy and Darwinism},
  pdfauthor={H. Høffding},
  hidelinks
}

\pagestyle{fancy}
\fancyhf{}
\rhead{Høffding, \emph{Philosophy and Darwinism} (1874)}
\cfoot{\thepage}

\setstretch{1.3}

\title{\textbf{Philosophy and Darwinism}}
\author{H.\ Høffding\\[4pt]
  \small Translated from the Danish}
\date{\emph{Nær og Fjern}, 1874}

\begin{document}
\maketitle
\thispagestyle{fancy}

% ------------------------------------------------------------------
%  NOTE: Part I is translated; Part II (the second half of the
%  original article) still requires translation.
%  Danish source: readings/hoffding-darwin/transcription.tex
% ------------------------------------------------------------------

\bigskip

\noindent It is a familiar fact of experience that oppositions set
reflection in motion. What one previously apprehended with
indifference and without attention suddenly becomes interesting once
one catches sight of the oppositions bound up with it, and of the
tensions to which it gives rise. Someone who, for other reasons, would
feel no inclination to object to Darwinism may nevertheless find a
motive for doing so in the conflicting interpretations of it---the
diametrically opposed judgments of its worth and significance that one
encounters on a daily basis. For one person, Darwinism is a wild
speculation, a product of our materialistic age, which seeks
confirmation for its strange and ultimate drives and inclinations in
self-generated hypotheses; for another, it is a brilliant and
magnificent conception, opening up vast perspectives, satisfying the
human mind's ideal longing for the infinite, and awakening enthusiasm
for higher and progressive forms of development. And if we return to
the narrower circle in which this doctrine arose, and where scientific
testing properly belongs---to the natural scientists themselves---we
likewise find a striking opposition of views. While a large number of
natural scientists (especially English and German ones) regard Darwin's
doctrine as so highly probable that they are inclined to overlook what
is lacking in the evidential support---indeed, that several even
embark on far-reaching speculations and abandon the firm ground of
their discipline---others (notably French researchers, with whom most
Danes seem to side) insist that an explanation is not a proof, but
presupposes proof in order to have validity, and that, despite all the
importance hypotheses may have for science, it is entirely
inadmissible, even for a moment, to forget the gaps in the evidence
because of the appeal of a hypothetical explanation.

\medskip

\noindent There is therefore good reason to undertake an examination of
the logical and ethical worth of Darwinism, in order to settle which
of these competing judgments is the true one, and what stance we
ought to adopt toward this remarkable doctrine, which has now made its
way from the scholar's quiet study out into the streets and is being
employed as a powerful weapon in the intellectual battles of our age.
Philosophy has, apart from these external occasions, by virtue of its
own task, a duty to carry out such an examination. It is philosophy's
task to attain self-consciousness with respect to the nature and method
of scientific cognition, to sharpen insight into the stages a doctrine
must pass through before it can be incorporated into the realm of
scientific truths, and to uphold the necessary conditions for the truth
and validity of cognition. But its vocation is not confined to this.
As it investigates the life of human consciousness, its laws and its
elements, and traces its development through history, it seeks at the
same time to situate every cognitive result as a member within the
whole organism of mind, and must therefore at every point examine how
the result achieved stands as a contribution to a general worldview.

\medskip

\noindent It is clear that the logical examination is here the most
important. One must first settle whether a doctrine is true or not,
before proceeding to examine its significance in a metaphysical and
ethical respect.

\section*{I.}

One may say that science would stand still without hypothesis. In
order to find the law according to which a given factor operates, or
to find the cause underlying a given regularity in phenomena, one must
necessarily venture beyond what is immediately at hand, and must often
make large leaps into entirely different circles of phenomena in order
to test whether the sought explanation is not to be found there.

\medskip

\noindent What is it, then, that Darwin's hypothesis seeks to explain?
It is the fact that living beings are divided into certain distinct
groups with distinctive characteristic marks that remain fixed for
each individual group---that is, the fact of what are called specific
differences.

\medskip

\noindent The resolution of the dispute depends on the question: does
natural selection have limits or not? But this question science is
unable to answer at its current stage. It is this that makes
Darwinism a hypothesis and justifies the hesitation of sober-minded
researchers to see attempts to present it as an established doctrine
that one may now simply accept.

\medskip

\noindent Darwin has demonstrated a real cause that is of essential
and perhaps hitherto overlooked significance for organic life, but he
has not been able to prove that this cause contains the sufficient
ground for the origin of the different species.

\medskip

\noindent That the species must have a natural origin, that the causes
of their arising, whatever they may otherwise be, must be sought in
the conditions given in nature---this is a presupposition that arises
from the principles of all scientific inquiry, principles that are
increasingly being applied even to domains from which they were
previously excluded.

\section*{II.}

% ------------------------------------------------------------------
%  TODO: Translate the second half of the article (Part II onward).
%  The Danish original begins: "Den logiske undersøgelse af
%  Darwinismen fører således til at anvise den dens rette Plads..."
% ------------------------------------------------------------------

\noindent[\emph{Translation of Part II forthcoming.}]

\end{document}
